%! Operations with Polynomials — Unit 3, Lesson 3.1 — Guided Notes
\documentclass[10pt]{article}
\usepackage{algebra2-article}
\usepackage{algebra2-boxes}
\newcommand{\TallMath}[1]{$\displaystyle #1\rule[-1.4em]{0pt}{3.2em}$}

\begin{document}

\pageheader{Unit 3, Lesson 3.1}{Guided Notes}
\namedateperiod

\vspace{0.08in}

\begin{objectivebox}
\textbf{By the end of this lesson, I will be able to\dots}
\begin{itemize}\setlength{\itemsep}{2pt}
  \item write a polynomial in \emph{standard form} and name its \emph{degree}, \emph{leading coefficient}, and number of \emph{terms};
  \item \emph{add} and \emph{subtract} polynomials by combining like terms, distributing a subtraction to \emph{every} term;
  \item \emph{multiply} polynomials in one and two variables and apply the \emph{special products} $(a\pm b)^2$ and $(a+b)(a-b)$;
  \item predict a product's \emph{degree} and \emph{leading coefficient} and state the \emph{end behavior} they force.
\end{itemize}
\end{objectivebox}

\vspace{0.05in}

\begin{vocabbox}
\small
Fill in each term as we build it together.
\vspace{2pt}
\termblanklong{Standard form}
\termblanklong{Like terms}
\termblanklong{Monomial / binomial / trinomial}
\termblanklong{Degree of a product}
\termblanklong{Special products}
\end{vocabbox}

\begin{hookbox}
Yesterday we used \emph{two} names for the same curve:
\[
  g(x) = (x-1)^2(x+2) \qquad\text{and}\qquad g(x) = x^3 - 3x + 2 .
\]
Check a few inputs in \emph{both} forms:
\begin{center}
\renewcommand{\arraystretch}{1.3}
\begin{tabular}{|c|c|c|}
\hline
$x$ & $(x-1)^2(x+2)$ & $x^3-3x+2$ \\ \hline
$0$ & \blank{1.4cm} & \blank{1.4cm} \\ \hline
$2$ & \blank{1.4cm} & \blank{1.4cm} \\ \hline
$-2$ & \blank{1.4cm} & \blank{1.4cm} \\ \hline
\end{tabular}
\end{center}
\begin{itemize}\setlength{\itemsep}{1pt}
  \item Agreeing at three points is good evidence. Is it \emph{proof} that the two forms are the same function? Why or why not? \writeline
  \item What would actually prove it? \writeline
\end{itemize}
Each costume is useful: the \textbf{factored} form hands you the \emph{zeros}; the \textbf{standard}
form hands you the \emph{degree} and \emph{leading coefficient}, hence the \emph{end behavior}. Today
we learn to travel from the first to the second.
\end{hookbox}

\begin{notesbox}{1. Naming a Polynomial}
Always put the polynomial in \textbf{standard form} first --- powers from \blank{2.6cm} to
\blank{2.6cm} --- because every other name depends on it.

\vspace{3pt}
Rewrite $-5x^2 + 7x^4 - 3 + x$ in standard form: \blank{6.5cm}

\vspace{3pt}
\begin{itemize}[leftmargin=1.5em, itemsep=2pt]
  \item \textbf{Degree:} \blank{1.0cm} \qquad \textbf{Leading coefficient:} \blank{1.0cm} \qquad
        \textbf{Constant term:} \blank{1.0cm}
  \item \textbf{Number of terms:} \blank{1.0cm}, so it is a (monomial / binomial / trinomial / none of
        these): \blank{2.6cm}
  \item In \emph{two} variables, the degree of a \emph{term} is the \textbf{sum} of its exponents. So
        $-7x^2y^3$ has degree \blank{1.0cm}.
\end{itemize}
\end{notesbox}

\begin{notesbox}{2. Adding \& Subtracting: Collecting, Not Creating}
\textbf{Like terms} have the \emph{same} variables raised to the \emph{same} exponents. So $3x^2$ and
$3x^3$ \blank{3.0cm} combine.

\vspace{4pt}
\textbf{Add.} $(3x^3 - 5x + 2) + (x^3 + 4x^2 - 7) = $ \blank{6.0cm}

\vspace{5pt}
\textbf{Subtract --- the whole lesson in one move.} A minus sign in front of a parenthesis is a
factor of $-1$, so it changes the sign of \blank{3.6cm} term inside.

\vspace{3pt}
$(3x^3 - 5x + 2) - (x^3 + 4x^2 - 7)$

\vspace{2pt}
\emph{Step 1 --- rewrite with every sign changed:} \; $3x^3 - 5x + 2$ \blank{5.0cm}

\vspace{2pt}
\emph{Step 2 --- collect like terms:} \blank{6.5cm}

\vspace{5pt}
\textbf{A degree surprise.} $(x^4 + 2x) - (x^4 - 5) = $ \blank{4.0cm}, which has degree
\blank{1.0cm} --- not $4$! The leading terms \blank{2.4cm}, so a sum or difference has degree
\emph{at most} the larger of the two degrees.
\end{notesbox}

\begin{notesbox}{3. Multiplying: Every Term Times Every Term}
One rule --- the \textbf{distributive property} --- in three sizes. Remember: coefficients
\blank{1.8cm}, exponents \blank{1.8cm}.

\vspace{4pt}
\textbf{(a) Monomial $\times$ polynomial.} \; $-2x^2(3x^3 - 4x + 6) = $ \blank{6.0cm}

\vspace{5pt}
\textbf{(b) Binomial $\times$ binomial} (four products --- ``FOIL'' is just a name for the same
distribution). \; $(2x - 5)(3x + 4)$

\vspace{2pt}
Four products: \blank{2.0cm} $+$ \blank{2.0cm} $+$ \blank{2.0cm} $+$ \blank{2.0cm}

\vspace{2pt}
Collected: \blank{5.0cm}

\vspace{5pt}
\textbf{(c) Binomial $\times$ trinomial} (six products --- use the box so nothing is dropped).
\; $(x + 3)(x^2 - 2x + 4)$
\begin{center}
\renewcommand{\arraystretch}{1.9}
\begin{tabular}{|c|c|c|c|}
\hline
$\times$ & $x^2$ & $-2x$ & $+4$ \\ \hline
$x$   & \blank{1.6cm} & \blank{1.6cm} & \blank{1.6cm} \\ \hline
$+3$  & \blank{1.6cm} & \blank{1.6cm} & \blank{1.6cm} \\ \hline
\end{tabular}
\end{center}
Collect like terms: \blank{6.5cm}

\vspace{5pt}
\textbf{(d) Two variables.} \; $(3x + 2y)(x - 4y) = $ \blank{6.0cm}

\vspace{2pt}
The two middle products, $-12xy$ and $+2xy$, are like terms because they have the
\blank{4.4cm}.

\vspace{5pt}
\textbf{(e) The shape of a product --- before you expand.} In $(2x^3 - 1)(-3x^4 + x)$, only the two
\emph{leading} terms can make the highest power: $2x^3 \cdot (-3x^4) = $ \blank{1.8cm}.
\begin{itemize}[leftmargin=1.5em, itemsep=2pt]
  \item Degrees \blank{1.4cm}: \; $3 + 4 = $ \blank{0.8cm}. \quad Leading coefficients
        \blank{1.8cm}: \; $2 \cdot (-3) = $ \blank{0.9cm}.
  \item So (Lesson 3.0) the degree is \emph{odd} and the leading coefficient is \emph{negative}:
        the left arm goes \blank{1.2cm} and the right arm goes \blank{1.2cm}.
\end{itemize}
\end{notesbox}

\begin{notesbox}{4. Special Products --- Worth Knowing on Sight}
Derive them once, then recognize them.
\begin{itemize}[leftmargin=1.5em, itemsep=3pt]
  \item \textbf{Square of a sum:} $(a+b)^2 = (a+b)(a+b) = $ \blank{4.2cm}
  \item \textbf{Square of a difference:} $(a-b)^2 = $ \blank{4.2cm}
  \item \textbf{Difference of squares:} $(a+b)(a-b) = $ \blank{3.0cm} \; --- the two middle terms
        \blank{2.2cm}.
\end{itemize}

\vspace{4pt}
\textbf{The trap.} $(x+4)^2$ is \emph{not} $x^2 + 16$. Test it at $x = 1$: \;
$(1+4)^2 = $ \blank{1.2cm} \; but \; $1^2 + 16 = $ \blank{1.2cm}. \; The missing piece is the
\textbf{middle term} \blank{1.6cm}. Correct answer: \blank{4.0cm}

\vspace{4pt}
Now use the identities:
\begin{itemize}[leftmargin=1.5em, itemsep=3pt]
  \item $(3x - 4)^2 = $ \blank{5.0cm}
  \item $(2x + 7)(2x - 7) = $ \blank{4.0cm}
  \item $(5x - 3y)^2 = $ \blank{5.4cm}
\end{itemize}
\end{notesbox}

\begin{practicebox}
\textbf{Back to the Hook.} Prove that $(x-1)^2(x+2)$ really is $x^3 - 3x + 2$.
\begin{enumerate}[leftmargin=1.7em, itemsep=6pt]
  \item First square the binomial: $(x-1)^2 = $ \blank{4.4cm} \; (special product --- don't forget the
        middle term).
  \item Now multiply that trinomial by $(x+2)$. Six products:
        \begin{center}
        \renewcommand{\arraystretch}{1.9}
        \begin{tabular}{|c|c|c|c|}
        \hline
        $\times$ & $x^2$ & $-2x$ & $+1$ \\ \hline
        $x$   & \blank{1.5cm} & \blank{1.5cm} & \blank{1.5cm} \\ \hline
        $+2$  & \blank{1.5cm} & \blank{1.5cm} & \blank{1.5cm} \\ \hline
        \end{tabular}
        \end{center}
  \item Collect like terms and write the result in standard form: \blank{5.5cm}
  \item Does it match $x^3 - 3x + 2$? \blank{1.6cm} \; Its degree is \blank{0.8cm} and its leading
        coefficient is \blank{0.8cm}, so the arms go \blank{3.4cm} --- exactly the graph you read in
        Lesson 3.0.
\end{enumerate}
\end{practicebox}

\end{document}
